\section{The Fathers did not ask}

A study of this kind is obliged to report what it does not find as carefully as
what it does. The finding of this section is negative and it is material: among
the ancient Christian and Hellenistic-Jewish witnesses reached here, nobody
raises the question of the hour at Genesis 15:5 --- and several of them come
close enough to it that the omission is informative rather than accidental.

The result is bounded, and the bounds are stated at the end. It is not a claim
that no Father ever mentioned the difficulty. It is a claim about a named set of
witnesses, each read at its own locus.

\subsection{Augustine comes within a sentence of it}

Augustine gives Genesis 15 two chapters of the sixteenth book of \textit{The
City of God}, and they are not perfunctory. Chapter 23 treats the promise;
chapter 24 treats the sacrifice, the ecstasy, and the fire. He handles the stars
directly, and what he says about them is exactly the thing that ought to have
raised the question:

\begin{studyframe}
\textit{\dots{} quia nec omnes eas uideri posse credendum est. Nam quanto quisque
acutius intuetur, tanto plures uidet. Unde et acerrime cernentibus aliquas
occultas esse merito existimatur\dots{} Postremo quicumque uniuersum stellarum
numerum conprehendisse et conscripsisse iactantur, sicut Aratus uel Eudoxus uel
si qui alii sunt, eos libri huius contemnit auctoritas.}\endnote{Augustine,
\textit{De civitate Dei} XVI.23, in the Hoffmann CSEL~40 Latin registered in
this repository. Marcus Dods' 1871 English, also registered here, renders the
passage: ``it is not to be believed that all of them can be seen. For the more
keenly one observes them, the more does he see. So that it is to be supposed
some remain concealed from the keenest observers\dots{} Finally, the authority
of this book condemns those like Aratus or Eudoxus, or any others who boast that
they have found out and written down the complete number of the stars.''}
\end{studyframe}

Augustine is asking whether the stars are countable. He decides they are not,
and adds the empirical observation that a sharper eye sees more of them, and
then reaches for two Greek astronomers to be refuted. He is, in other words,
thinking about the visibility of stars in the very sentence where the visibility
of stars is the issue --- and the thought that they might not have been visible
at that hour does not arise.

Nor does it arise in chapter 24, where he quotes the whole of verses 7 to 21 and
comments on the temporal markers. He allegorizes them: ``that about the going
down of the sun great fear fell upon Abraham and a horror of great darkness,
signifies that about the end of this world believers shall be in great
perturbation and tribulation''; and the fire between the pieces ``signifies that
at the end of the world the carnal shall be judged by fire''. And he closes the
whole with the sentence that explains everything: \textit{Haec omnia in uisu
facta diuinitus adque dicta sunt} --- gloss: \emph{all these things were divinely
done and said in a vision}.\endnote{Augustine, \textit{De civitate Dei} XVI.24,
CSEL~40; the English clauses are Dods'. Augustine's Old Latin text of the chapter,
quoted in full at that place, has \textit{et consedit illis Abram} at verse 11
and \textit{flamma facta est} without any darkness at verse 17, agreeing with
the Greek.}

That is the whole mechanism. Augustine's Genesis 15 is a vision, its sunset is
the end of the world, and its clock is not a clock. He had no reason to ask what
time it was, because in his reading nothing in the chapter happens at a time.

\subsection{Philo has already dissolved the question}

The same structure, three centuries earlier and more thoroughly. Philo's
treatise on this chapter reads the going-out as the soul's exit from the body
and the senses; the setting sun as the setting of the human intellect; the
trance as prophetic possession. There is no place in that reading where an hour
could be asked about.

And Philo, like Augustine, comes close. Immediately after quoting the command to
count the stars he writes: ``But if it is impossible to find out the number of
the stars which are perceptible by the outward senses, how much more so must it
be to count those which are discernible only by the intellect?'' The
countability problem is raised; the visibility problem is not.\endnote{Philo,
\textit{Quis rerum divinarum heres sit}, chapter XVII, Yonge. In the same
chapter Philo describes the seed as ``like the ethereal firmament which thou
beholdest\dots{} full of unshadowed and pure brilliancy (for night is driven
away from heaven, and darkness from virtue)'' --- a sky, notably, from which
night has been removed.}

\subsection{Josephus removes the evidence entirely}

Josephus retells the chapter in the \textit{Antiquities}, and his paraphrase is
worth quoting in full because of what is missing from it:

\begin{studyframe}
And God promised, that he should have a son, and that his posterity should be
very numerous; insomuch, that their number should be like the stars. When he
heard that, he offered a sacrifice to God, as he commanded him.\endnote{Josephus,
\textit{Antiquitates Judaicae} I.7.2, in William Whiston's public-domain
translation, read in the Penelope (University of Chicago) transcription. The
passage continues with the animals, the birds of prey, and the four hundred
years, and never mentions a sunset, a darkness, or a day.}
\end{studyframe}

There is no bringing-out, no command to look, no counting, no sunset, no
darkness, and no ``that day''. Josephus has converted a demonstration into an
assertion and deleted the chapter's entire clock. Whatever his reasons, the
effect is that the earliest continuous retelling of Genesis 15 available to this
study contains none of the elements from which the difficulty is built.

\subsection{The Targums keep the sequence and add nothing}

This is the most probative of the negative results, because the Aramaic
translators were not shy.

Targum Onkelos renders the three markers closely: \textit{vahavah shimsha
lemei'al} at verse 12 (the sun about to enter), \textit{vahavah shimsha allat}
at verse 17 (the sun entered), with the darkness present. Targum
Pseudo-Jonathan is wildly expansive at verse 12, converting the dread and
darkness into the four kingdoms of Babylon, Media, Greece, and Persia; and at
verse 17 it has Abram see Gehinnom ascending with flaming coals.\endnote{Targum
Onkelos and Targum Pseudo-Jonathan on Gen 15:5, 12, 17, read in the Aramaic and,
for Pseudo-Jonathan, in J.~W.~Etheridge's public-domain English of 1862.
Etheridge renders Pseudo-Jonathan at 15:12: ``And when the sun was nearing to
set, a deep sleep was thrown upon Abram: and, behold, four kingdoms arose to
enslave his children''; and at 15:17: ``And when the sun had set there was
darkness.''}

A translator who will insert four world empires into verse 12 is a translator
who inserts what he thinks the text needs. Neither Targum inserts a night before
verse 5, a morning after verse 6, or any interval whatever between the promise
and the rite. Both keep the sequence exactly as the Hebrew has it.

\subsection{The later Catholic commentators notice everything except this}

Cornelius a Lapide is the test case, because he does raise the hour --- and then
does nothing with it. His note on verse 5 opens:

\begin{studyframe}
\textit{Numera stellas: erat ergo tunc nox, non illunis, sed innubis, serena et
stellata.}\endnote{A Lapide on Gen 15:5, 1700 printing, orthography normalized.
He immediately continues with the countability argument and supports it from
Augustine, reproducing the sentence about the keener eye with a small addition of
his own (\textit{tanto plures in coelo videt}) and citing \textit{De civitate
Dei} book 16; the chapter number in the copy read here is not legible with
confidence and is not reported. He
adds that God bids Abram number the stars ``both because he was an astronomer,
and because he was accustomed to look up at them'', and cites the verses
attributed to Orpheus in the fifth book of Clement's \textit{Stromateis} in
which Abraham is called an astrologer.}
\end{studyframe}

\noindent Gloss: \emph{Number the stars: it was therefore then night --- not
moonless, but cloudless, clear and starry.}

\textit{Erat ergo tunc nox.} The inference is exactly the one the daylight claim
denies, and a Lapide draws it in four words, from the command alone, without
consulting verse 12. Seven verses later he explains Abram's deep sleep by
``excessive daytime labour''. He has therefore committed himself to a night, a
following day of work, and an evening --- a two-day chronology --- inside a
chapter that ends ``in that day'', and he never says so. The two notes stand several columns
apart under separate verse headings, and nothing in his commentary brings them
into contact.

The others are simply silent. The 1609 Douay Old Testament annotates verses 6,
9, and 13 on this opening and prints nothing at verse 5. Challoner's revision
prints no note on the chapter at all beyond its one-line argument. Haydock
annotates verses 1, 2, 6, 8, 9, 12, 13, 14, and 16 --- and skips verse 5
entirely, having already used it, at verse 1, as his proof that Abram was awake.

\subsection{The bound of this negative result}

\begin{studybox}{What the negative result covers, and what it does not}
\textbf{Covered, each read at its own locus:} Philo, \textit{Quis rerum
divinarum heres sit} (chapters XVI--XVII, XX, LI--LIII); Josephus,
\textit{Antiquities} I.7.2; Targum Onkelos and Targum Pseudo-Jonathan on Genesis
15:5, 12, 17; Augustine, \textit{De civitate Dei} XVI.23--24 in both the CSEL
Latin and the Dods English; the 1609 Douay Old Testament apparatus on Genesis
15; Cornelius a Lapide on Genesis 15:1, 5, 10, 11, 12, and 17; Haydock's
annotations on Genesis 15. None of these raises the conflict between the
star-promise and the chapter's own sunset.

\textbf{Not covered, and named so that no clause of this study depends on
them:} Jerome's \textit{Hebrew Questions on Genesis}; Ambrose, \textit{De
Abraham}; John Chrysostom's homilies on Genesis; Theodoret's
\textit{Quaestiones in Genesim}; Ephrem; the Glossa ordinaria; Rupert of Deutz;
Pererius, Tostatus, and Oleaster. Ambrose and Theodoret were reached only
through citations in a Lapide and in the 1609 Douay margin, and are reported here
only as those authors report them. Jubilees, the Genesis Apocryphon, and Targum
Neofiti were not reached at all.

\textbf{What the covered set does establish:} that the difficulty was invisible
to the ancient Christian tradition for a specific and identifiable reason, and
not by oversight. The Greek text removed the darkness of verse 17; the Old Latin
made from it did the same; and the dominant Christian reading of the chapter ---
Philo's, and Augustine's after him --- converted every temporal marker in it into
an allegory of the soul or of the end of the world. A clock that has been turned
into a symbol cannot be read for the time.
\end{studybox}
